% =====================================================================
% Topic 4 — Number Sense — part a: theory, worked examples, problems 1–15
% =====================================================================
\begin{knowbox}
\textbf{Big idea.} ``Number sense'' is the ability to \emph{understand} numbers, not just to \emph{compute} with them:
knowing which number is larger, whether a result is reasonable, how to round for a quick estimate, and seeing the
hidden structure inside a number (its digits, its divisors, its multiples). In this topic you will review and connect:
\begin{itemize}[leftmargin=1.4em,itemsep=1pt,topsep=2pt]
  \item \textbf{decimals and negative numbers}: comparing and ordering them on the number line;
  \item the \textbf{order of operations} and \textbf{estimating} results;
  \item \textbf{fractions, ratios, percentages}: ``a fraction of a number'', equivalent fractions, discounts;
  \item \textbf{speed} -- distance -- time and unit conversion (km/h, m/s);
  \item \textbf{divisibility, common divisors, last digits, place value} (solving digit puzzles);
  \item the \textbf{mean, the median} and \textbf{number patterns} (Fibonacci, arithmetic sequences).
\end{itemize}
The general secret: \emph{estimate before you calculate; check after you calculate}. A student with good number
sense can often eliminate several unreasonable options in a multiple-choice question without any exact calculation.
\end{knowbox}
\vspace{0.3cm}

% ---------------------------------------------------------------------
\shead{Decimals: reading, comparing, ordering, adding and subtracting}
\begin{knowbox}
\textbf{1. Place value.} In a decimal, each digit sits in a \emph{place}. After the decimal point come the tenths,
the hundredths, the thousandths, \ldots
\begin{center}
\renewcommand{\arraystretch}{1.25}
\begin{tabular}{|c|c|c|c|c|c|}
\hline
\rowcolor{tmBlue!10} Tens & Ones & . & Tenths & Hundredths & Thousandths\\ \hline
 & $2$ & . & $2$ & $3$ & \\ \hline
 & $2$ & . & $3$ & $0$ & \\ \hline
\end{tabular}
\end{center}
Example: $2.23 = 2 + \dfrac{2}{10} + \dfrac{3}{100}$, while $2.3 = 2 + \dfrac{3}{10} = 2.30$.

\textbf{2. Adding or removing zeros at the end of the decimal part does not change the value:}
$2.3 = 2.30 = 2.300$. Reason: $\dfrac{3}{10} = \dfrac{30}{100} = \dfrac{300}{1000}$ (multiply numerator and denominator by $10$).

\textbf{3. Rule for comparing two positive decimals.}
\begin{itemize}[leftmargin=1.4em,itemsep=1pt,topsep=2pt]
  \item Compare the \emph{whole-number parts} first: the one with the larger whole part is larger.
  \item If the whole parts are equal, \emph{add zeros} so that both decimal parts have the same number of digits,
  then compare place by place from left to right: tenths, then hundredths, \ldots
\end{itemize}
Example: compare $0.304$ and $0.34$. Write $0.34 = 0.340$. Tenths: $3 = 3$; hundredths: $0 < 4$.
So $0.304 < 0.340 = 0.34$.

\textbf{4. Adding and subtracting decimals:} write the numbers in a column \emph{with the decimal points lined up}
(places lined up), then add/subtract as with whole numbers.

\textbf{5. Rounding decimals.} To round to the nearest hundredth (two digits after the point), look at the digit just
to the right of that place (the thousandths digit): if it is less than $5$, keep the hundredths digit; if it is $5$ or more,
add $1$ to the hundredths digit; then drop all digits after it. Example: $4.6875 \approx 4.69$ (the thousandths digit is
$7 \ge 5$), while $3.1416 \approx 3.14$ (because $1 < 5$). The sign ``$\approx$'' reads ``is approximately''.
\end{knowbox}
\begin{proofbox}
Why can't we compare ``by length'' (the number with more digits is larger)? Write both in thousandths:
$0.304 = \dfrac{304}{1000}$ and $0.34 = \dfrac{340}{1000}$. The fractions have the same denominator and $304 < 340$, so
$0.304 < 0.34$, even though $0.304$ ``looks longer''. Adding zeros is exactly bringing both to the same denominator.
\end{proofbox}
\begin{warnbox}
Many students think $2.23 > 2.3$ because ``$23 > 3$''. Wrong! Compare place by place: $2.23$ has $2$ tenths,
while $2.3$ has $3$ tenths, so $2.23 < 2.3$.
\end{warnbox}
\begin{tipbox}
To find the smallest (largest) number in a set of decimals, write all of them in \emph{one column} with the decimal
points lined up, adding zeros to make them equally long. Then you just compare them like whole numbers.
\end{tipbox}

% ---------------------------------------------------------------------
\shead{Negative numbers and the number line}
\begin{knowbox}
\textbf{1. The number line.} The integers are placed on a straight line: $0$ in the middle, positive numbers to the
right, negative numbers to the left.
\numlinex{8}{0/$-4$,1/$-3$,2/$-2$,3/$-1$,4/$0$,5/$1$,6/$2$,7/$3$,8/$4$}{}{4}
\textbf{2. Comparison rule.} On a (horizontal) number line, \emph{a number to the left is smaller than a number to the
right}. Therefore:
\begin{itemize}[leftmargin=1.4em,itemsep=1pt,topsep=2pt]
  \item every negative number is less than $0$, and $0$ is less than every positive number;
  \item every negative number is less than every positive number;
  \item of two negative numbers, the one \emph{farther from $0$} is \emph{smaller}: $-8 < -2$ because $8 > 2$.
\end{itemize}
\textbf{3. Everyday pictures.} A temperature of $-8^\circ$C is colder than $-2^\circ$C; owing $8$ dollars
leaves you ``poorer'' than owing $2$ dollars (an amount of $-8$ compared with $-2$).

\textbf{4. The coordinate plane.} A point is given by a pair of numbers $(x;y)$: $x$ is the horizontal coordinate
(moving right/left), $y$ is the vertical coordinate (moving up/down). A horizontal move changes only $x$; a vertical
move changes only $y$.
\end{knowbox}
\begin{proofbox}
Why is $-8 < -2$? We have $-2 = -8 + 6$: to go from $-8$ to $-2$ on the number line we must move $6$ units \emph{to the
right}. Moving right means increasing, so $-2$ is larger than $-8$.
\end{proofbox}
\begin{warnbox}
Do not compare negative numbers ``by the feel of positive numbers''. $-8$ is \emph{less} than $-2$, even though $8$ is
greater than $2$. When in doubt, sketch a quick number line.
\end{warnbox}
\begin{tipbox}
\textbf{Parity (even -- odd)} is a powerful ``number sense'' tool: even $+$ even $=$ even, odd $+$ odd $=$ even,
even $+$ odd $=$ odd; a product of odd numbers is always odd. For instance, the sum of $5$ odd numbers is odd, and the sum
of $4$ odd numbers is even. When a point moves on a coordinate grid, looking at the parity of its coordinates often tells
you at once whether a position can be reached.
\end{tipbox}

% ---------------------------------------------------------------------
\shead{Order of operations and estimation}
\begin{knowbox}
\textbf{1. Order of operations.}
\begin{enumerate}[leftmargin=1.6em,itemsep=1pt,topsep=2pt]
  \item \emph{Brackets} first: $(\;)$, then $[\;]$, then $\{\;\}$.
  \item Next, \emph{powers}.
  \item Then \emph{multiplication and division}, from left to right.
  \item Finally \emph{addition and subtraction}, from left to right.
\end{enumerate}
Example: $6\times 2 - 3 = 12 - 3 = 9$ (multiply first, then subtract), whereas $6\times(2-3) = 6\times(-1) = -6$.

\textbf{2. Rounding to estimate.} To know quickly ``roughly how much'' a result is, round the factors to the nearest
ten or hundred and compute mentally. Example: $99 \times 9 \approx 100 \times 9 = 900$.

\textbf{3. A quick-calculation trick:} $99 \times 9 = (100-1)\times 9 = 900 - 9 = 891$ (multiplication distributes
over subtraction).
\end{knowbox}
\begin{proofbox}
Why is $(100-1)\times 9 = 100\times 9 - 1\times 9$? Imagine $9$ rows of $100$ sweets each, but $1$ sweet is taken from
every row. The sweets left are $9$ rows $\times$ $99$ sweets, which also equals ($9\times 100$ sweets) minus (the $9$
sweets taken). This is the \emph{distributive property}: $a\times(b-c) = a\times b - a\times c$.
\end{proofbox}
\begin{warnbox}
Working ``from left to right regardless of the operation'' is wrong: $2 + 3\times 4 \ne 5\times 4$. The correct
value is $2 + 3\times 4 = 2 + 12 = 14$.
\end{warnbox}
\begin{tipbox}
In a multiple-choice question asking ``which number is closest to\ldots'', the options often differ by factors of $10$
($10$; $100$; $1000$; \ldots). Then a rough estimate is enough to choose correctly.
\end{tipbox}

% ---------------------------------------------------------------------
\shead{Fractions: a fraction of a number, equivalent fractions}
\begin{knowbox}
\textbf{1. A fraction is a part of a whole.} $\dfrac{a}{b}$ means: split the whole into $b$ equal parts and take $a$ of them.

\textbf{2. Finding a fraction of a number:} to find $\dfrac{m}{n}$ of $a$, compute $a \times \dfrac{m}{n}$
(split $a$ into $n$ parts and take $m$ of them). Example: $\dfrac{1}{3}$ of $396$ is $396 \div 3 = 132$.

\textbf{3. The part that is left.} If $\dfrac{m}{n}$ of a job is done, the part left is $1 - \dfrac{m}{n} = \dfrac{n-m}{n}$.

\textbf{4. Equivalent fractions (the basic property).} Multiplying (or dividing) the numerator and the denominator by the
same non-zero number gives an equal fraction:
\[ \frac{3}{7} = \frac{3\times 9}{7\times 9} = \frac{27}{63}, \qquad \frac{12}{36} = \frac{12\div 12}{36\div 12} = \frac{1}{3}. \]
\textbf{5. Simplifying a fraction:} divide the numerator and the denominator by their greatest common divisor to get a
fraction in lowest terms.

\textbf{6. Two numbers with a known ratio and sum (or difference).} If $\dfrac{x}{y} = \dfrac{5}{8}$, then there is a
number $k$ with $x = 5k$ and $y = 8k$ (because $\dfrac{5k}{8k} = \dfrac58$). Knowing the sum $x + y$ (or the difference
$y - x$), we form an equation in $k$ and solve it. Example: $x + y = 91$ gives $5k + 8k = 91$, i.e.\ $13k = 91$, $k = 7$;
hence $x = 35$, $y = 56$.

\textbf{7. A fraction of a fraction.} ``Half of a quarter'' is $\dfrac12\times\dfrac14 = \dfrac18$: multiply numerator by
numerator, denominator by denominator. This happens often when a shape is halved several times in a row: each halving
multiplies the area fraction by $\dfrac12$ again.

\textbf{8. Comparing fractions.} Bring them to a common denominator and compare numerators, or convert them to
decimals. Examples: $\dfrac23 = \dfrac{8}{12} > \dfrac{7}{12}$; \ $\dfrac35 = 0.6 < \dfrac34 = 0.75$.
\end{knowbox}
\begin{proofbox}
Why does multiplying the numerator and denominator by the same number leave the fraction unchanged? $\dfrac{3}{7}$ is
``cut the cake into $7$ parts, take $3$''. If we cut \emph{each} part into $9$ smaller pieces, the cake has
$7\times 9 = 63$ pieces and our share has $3\times 9 = 27$ pieces. The amount of cake has not changed, so
$\dfrac{3}{7} = \dfrac{27}{63}$.
\end{proofbox}
\begin{warnbox}
Read the question carefully: does it ask for the part \emph{done} or the part \emph{left}? Many students correctly find
$\dfrac13$ of $396$ to be $132$ and circle it straight away, while the question asks for the pages \emph{left},
$396 - 132 = 264$.
\end{warnbox}
\begin{parentbox}
In Grade 7, encourage your child to \emph{use an unknown}: let $x$ stand for the quantity we do not know, write the
conditions of the problem as an equation, solve it and check. For example, ``read $\dfrac13$ of a $396$-page book, $x$
pages left'' becomes $132 + x = 396$. This works for every problem, even when the numbers are ``awkward'' and a picture
would not help.
\end{parentbox}

% ---------------------------------------------------------------------
\shead{Percentages, unit price and discounts}
\begin{knowbox}
\textbf{1. Percent means ``out of one hundred''.} $80\% = \dfrac{80}{100} = 0.8$.

\textbf{2. Finding $p\%$ of a number $a$:} compute $a \times \dfrac{p}{100}$. Example: $80\%$ of $50$ is
$50 \times \dfrac{80}{100} = 40$.

\textbf{3. Expressing $a$ as a percentage of $b$:} compute $\dfrac{a}{b}\times 100\%$. Example: $32$ correct answers out of
$40$ questions is $\dfrac{32}{40}\times 100\% = 80\%$.

\textbf{4. Unit price.} The unit price is the price of \emph{one} unit (one kilogram, one book, \ldots):
\[ \text{Total cost} = \text{unit price} \times \text{quantity}, \qquad \text{unit price} = \text{total cost} \div \text{quantity}. \]
\textbf{5. A discount of $p\%$.} New price $=$ old price $\times \left(1 - \dfrac{p}{100}\right)$. A $10\%$ discount
leaves $90\%$: multiply by $0.9$.

\textbf{6. Successive discounts.} Each time the price is reduced by $10\%$ \emph{of the current price}, we multiply by
$0.9$ once more. Example: $100$ dollars after two reductions: $100 \to 90 \to 81$ (dollars), \emph{not} $80$.
\end{knowbox}
\begin{proofbox}
Why are two $10\%$ reductions not the same as one $20\%$ reduction? The second reduction is $10\%$ of $90$ (that is $9$),
not $10\%$ of $100$ (that is $10$). Since the price has already dropped, the next reduction is smaller too: in total the
price drops by only $10 + 9 = 19$ dollars, i.e.\ $19\%$.
\end{proofbox}
\begin{tipbox}
Finding $10\%$ very fast: just divide by $10$ (move the decimal point one place to the left). From $10\%$ you get
$5\%$ (halve it), $20\%$ (double it), $80\% = 100\% - 20\%$, \ldots
\end{tipbox}

% ---------------------------------------------------------------------
\shead{Speed, distance, time and unit conversion}
\begin{knowbox}
\textbf{1. Three linked formulas.} For motion at a constant speed:
\[ s = v \times t, \qquad v = s \div t, \qquad t = s \div v, \]
where $s$ is the distance, $v$ the speed and $t$ the time.

\textbf{2. Who is faster?} The faster person covers a \emph{longer} distance in the \emph{same amount of time}. The
safest way: compute each person's speed in the same unit (for example m/s) and compare.

\textbf{3. Running laps.} Running $n$ laps of a track $L$ metres long means a distance of $n \times L$ metres.

\textbf{4. Converting km/h to m/s.} $1$ km $= 1000$ m and $1$ hour $= 3600$ seconds, so
\[ 1 \text{ km/h} = \frac{1000 \text{ m}}{3600 \text{ s}} = \frac{5}{18} \text{ m/s}, \qquad
   18 \text{ km/h} = 18\times\frac{5}{18} = 5 \text{ m/s}, \qquad 1 \text{ m/s} = 3.6 \text{ km/h}. \]
\end{knowbox}
\begin{proofbox}
Why $v = s \div t$ and $t = s \div v$? It is just finding an unknown factor in the equation $s = v\times t$: dividing
both sides by $t$ (non-zero) gives $v = \dfrac{s}{t}$; dividing both sides by $v$ (non-zero) gives $t = \dfrac{s}{v}$.
\end{proofbox}
\begin{warnbox}
The quantities must use \emph{matching units}: if the distance is in metres and the time in seconds, the speed comes
out in m/s; do not divide metres by hours. Before comparing two speeds, convert them to the same unit.
\end{warnbox}

% ---------------------------------------------------------------------
\shead{Divisibility, common divisors, last digits and place value}
\begin{knowbox}
\textbf{1. Multiples and divisors.} $a$ is a \emph{multiple} of $b$ (or $b$ is a \emph{divisor} of $a$) if
$a = b\times k$ for some whole number $k$. Example: $45 = 9\times 5$, so $45$ is a multiple of $9$.

\textbf{2. Divisibility rules.}
\begin{itemize}[leftmargin=1.4em,itemsep=1pt,topsep=2pt]
  \item by $2$: the last digit is $0, 2, 4, 6, 8$; \quad by $5$: the last digit is $0$ or $5$; \quad by $10$: the last digit is $0$;
  \item by $3$: the digit sum is divisible by $3$; \quad by $9$: \emph{the digit sum is divisible by $9$}.
\end{itemize}
\textbf{3. Common divisors.} A number that divides both $a$ and $b$ is a common divisor of $a$ and $b$. Every common divisor
is a divisor of the \emph{greatest common divisor} $\gcd(a,b)$. Example: $\gcd(40,72)=8$, so the common divisors of $40$
and $72$ are $1, 2, 4, 8$.

\textbf{4. The last digit of a product} depends only on the last digits of the factors:
$\ldots 3 \times \ldots 7$ ends in the same digit as $3\times 7 = 21$, namely $1$.

\textbf{5. Place value.} A number written with digits can be expanded by place:
\[ \nd{abc} = 100a + 10b + c, \qquad \nd{PQQ} = 100P + 10Q + Q = 100P + 11Q. \]
(The bar on top shows that this is the number with digits $a, b, c$, not the product $a\times b\times c$.)
\end{knowbox}
\begin{proofbox}
Why does the divisibility rule for $9$ work? Take the number $\nd{abc} = 100a + 10b + c$ and split it:
\[ 100a + 10b + c = (99a + 9b) + (a + b + c) = 9\times(11a + b) + (a+b+c). \]
The part $9\times(11a+b)$ is always divisible by $9$. So $\nd{abc}$ is divisible by $9$ exactly when the remaining part
$a+b+c$ (the digit sum) is divisible by $9$. For numbers with more digits the argument is the same, because
$10 = 9+1$, $100 = 99+1$, $1000 = 999+1$, \ldots
\end{proofbox}
\begin{tipbox}
For a ``digit puzzle'' (each letter stands for a digit) such as $\nd{PQQ}+\nd{PPQ}+\nd{QQQ}=876$, work \emph{from the ones
column to the left}, watching out for carries. The ones column usually has the fewest possibilities, so the first digit
is found very quickly. You can also use place value to turn the puzzle into one equation. When you have to count or find
all possibilities (sets of digits, outcomes of tossing coins, \ldots), \emph{list them systematically} (in increasing order,
or with a table/tree diagram) so that nothing is missed or repeated.
\end{tipbox}
\begin{knowbox}
\textbf{Systematic listing and chance.}
\begin{itemize}[leftmargin=1.4em,itemsep=1pt,topsep=2pt]
  \item To count the three-digit numbers whose digits have a given \emph{product}, find the sets of three digits (written in
  increasing order so as not to repeat), then count the arrangements of each set: three different digits can be arranged in
  $3\times 2\times 1 = 6$ ways; two equal digits and one other in $3$ ways. Example: product $12$ with the set $\{1;3;4\}$
  gives $6$ numbers $134, 143, 314, 341, 413, 431$; the set $\{2;2;3\}$ gives $3$ numbers $223, 232, 322$.
  \item When all outcomes are \emph{equally likely}, the probability of an event is
  \[ \frac{\text{number of favourable outcomes}}{\text{total number of possible outcomes}}. \]
  Example: tossing $2$ coins has $4$ equally likely outcomes: HH, HT, TH, TT (H: heads, T: tails). The probability that both
  coins show the same face is $\dfrac{2}{4} = \dfrac12$. Use a tree diagram to list them:
\end{itemize}
\begin{figC}
  \node[draw=Navy,circle,inner sep=1.5pt] (r) at (0,0) {\scriptsize start};
  \node[draw=tmBlue,fill=tmBlue!10,circle,inner sep=2pt] (a) at (2,0.8) {H};
  \node[draw=tmOrange,fill=tmOrange!10,circle,inner sep=2pt] (b) at (2,-0.8) {T};
  \foreach \n/\y/\l/\c in {aa/1.2/H/tmBlue,ab/0.4/T/tmOrange}{\node[draw=\c,fill=\c!10,circle,inner sep=2pt] (\n) at (4,\y) {\l};}
  \foreach \n/\y/\l/\c in {ba/-0.4/H/tmBlue,bb/-1.2/T/tmOrange}{\node[draw=\c,fill=\c!10,circle,inner sep=2pt] (\n) at (4,\y) {\l};}
  \draw[Slate] (r)--(a) (r)--(b) (a)--(aa) (a)--(ab) (b)--(ba) (b)--(bb);
  \node[right,font=\scriptsize] at (4.3,1.2) {HH \ (same face)};
  \node[right,font=\scriptsize] at (4.3,0.4) {HT};
  \node[right,font=\scriptsize] at (4.3,-0.4) {TH};
  \node[right,font=\scriptsize] at (4.3,-1.2) {TT \ (same face)};
\end{figC}
\end{knowbox}
\begin{warnbox}
``The digit sum is divisible by $9$'' is the rule for $9$; it does \emph{not} work for $4$, $6$ or $8$. And a number
divisible by $3$ is not necessarily divisible by $9$: $12$ is divisible by $3$ but not by $9$.
\end{warnbox}

% ---------------------------------------------------------------------
\shead{Number patterns, the mean and the median}
\begin{knowbox}
\textbf{1. The Fibonacci sequence.} $1, 1, 2, 3, 5, 8, 13, 21, 34, 55, \ldots$ Each number (from the third on) is
\emph{the sum of the two numbers just before it}: $2 = 1+1$, $3 = 1+2$, $5 = 2+3$, \ldots

\textbf{2. Arithmetic sequences (equally spaced).} Each number (from the second on) equals the previous number plus a fixed
number $d$ (the common difference). Example: $1, 4, 7, 10$ (with $d = 3$). The $n$-th number is $a_1 + (n-1)\times d$.

\textbf{3. The mean (average).} The mean of $n$ numbers $=$ (their sum) $\div n$. Hence
\[ \text{Sum} = \text{mean} \times \text{how many numbers}. \]
\textbf{4. The ``deviation'' method.} The deviations from the mean always add up to $0$: the ``excess'' of the numbers
above the mean exactly equals the ``shortfall'' of the numbers below the mean.

\textbf{5. The median.} Arrange the numbers in increasing order; the number in the middle is the median (when there is
an odd number of numbers). Example: the median of $42, 43, 46, 49, 50$ is $46$. With an even number of numbers, the median
is the mean of the two middle numbers: the median of $6, 7, 9, 10$ is $(7+9)\div 2 = 8$. Note: when a number is added or
changed, the ``middle'' position can move, so re-order the list.

\textbf{6. Consecutive integers.} For an odd number of consecutive integers, the mean is the middle number:
the mean of $19, 20, 21, 22, 23$ is $21$.
\end{knowbox}
\begin{proofbox}
Why do the deviations add up to $0$? Suppose $4$ numbers have mean $80$, so their sum is $4\times 80 = 320$. Write
each number as $80 + $ (its deviation). Adding up: $320 + $ (sum of deviations) $= 320$, so the sum of the deviations
must be $0$.
\end{proofbox}
\begin{figC}
  \draw[Slate!60] (-0.3,0)--(5.3,0);
  \node[left,font=\scriptsize,text=Slate] at (-0.3,0) {mean $=40$};
  \foreach \x/\v/\c in {0/38/tmOrange,1.5/43/tmGreen,3/41/tmGreen,4.5/38/tmOrange}{
    \pgfmathsetmacro\h{(\v-40)*0.3}
    \fill[\c!70] (\x,0) rectangle (\x+0.7,\h);
    \node[font=\scriptsize] at (\x+0.35,{\h+0.25*sign(\h+0.001)}) {$\v$};}
\end{figC}
\begin{center}\small The four numbers $38, 43, 41, 38$ have mean $40$. Deviations from $40$: $-2,\ +3,\ +1,\ -2$; sum $= 0$.\end{center}
\begin{parentbox}
You can illustrate the mean with stacks of coins: ``level out'' the tall stacks into the short ones until all stacks are
the same height. That common height is the mean. The coins moved out of the tall stacks exactly equal the coins added to
the short ones -- this is the idea ``the deviations add up to $0$''.
\end{parentbox}

% =====================================================================
% WORKED EXAMPLES
% =====================================================================
\needspace{14\baselineskip}
\par\vspace{0.7em}\noindent{\large\bfseries\color{tmOrange!90!black}\faPencil*\ \ Worked examples}\par\vspace{0.3em}

\begin{exbox}{Example 1}
Consider the set $\{0.5;\ 0.45;\ 0.405;\ 0.054;\ 0.54\}$.
\begin{enumerate}[label=\alph*),leftmargin=1.6em,itemsep=0pt,topsep=2pt]
  \item Arrange the numbers in increasing order.
  \item Find the sum of the largest and the smallest number.
\end{enumerate}
\guide
\step{1}{Write the numbers in a column, adding zeros to get three digits after the point}
Adding zeros at the end of the decimal part does not change the value, so:
\[ 0.5 = 0.500;\quad 0.45 = 0.450;\quad 0.405;\quad 0.054;\quad 0.54 = 0.540. \]
\step{2}{Compare them like whole numbers}
Dropping ``$0.$'', we compare $500;\ 450;\ 405;\ 54;\ 540$:
\[ 54 < 405 < 450 < 500 < 540. \]
So the increasing order is: $0.054 < 0.405 < 0.45 < 0.5 < 0.54$.
\step{3}{Add the largest and the smallest (decimal points lined up)}
\[ \begin{array}{r} 0.540\\ +\ 0.054\\ \hline 0.594 \end{array} \]
\verify{$0.594 - 0.54 = 0.054$. Correct.}
\ans{a) $0.054 < 0.405 < 0.45 < 0.5 < 0.54$;\quad b) $0.594$}
\end{exbox}

\begin{exbox}{Example 2}
Emma cycles along a $45$ km route. In the morning she covers $\dfrac{2}{5}$ of the route. How many kilometres does Emma
still have to ride in the afternoon, and what fraction of the whole route is that?
\guide
\step{1}{Distance covered in the morning}
\begin{recallbox}To find $\dfrac{m}{n}$ of a number $a$, compute $a\times\dfrac{m}{n}$.\end{recallbox}
Morning: $45 \times \dfrac{2}{5} = \dfrac{45\times 2}{5} = \dfrac{90}{5} = 18$ (km).
\step{2}{Introduce an unknown and form an equation}
Let $x$ (km) be the distance still to ride in the afternoon ($x > 0$). The morning distance plus the afternoon distance is
the whole route:
\[ 18 + x = 45. \]
\step{3}{Solve the equation}
$x = 45 - 18 = 27$ (km).
\step{4}{The fraction left}
Method 1: $1 - \dfrac{2}{5} = \dfrac{5}{5} - \dfrac{2}{5} = \dfrac{3}{5}$.\quad
Method 2: $\dfrac{27}{45} = \dfrac{27\div 9}{45\div 9} = \dfrac{3}{5}$. Both methods give the same result.
\verify{$18 + 27 = 45$ km and $45\times\dfrac35 = 27$ km. Correct.}
\ans{$27$ km, which is $\dfrac{3}{5}$ of the route}
\end{exbox}

\begin{exbox}{Example 3}
a) Evaluate $20 - 4\times 3 + 8 \div 2$ \quad and \quad $(20-4)\times(3+8) \div 2$.\\
b) Without calculating exactly, decide which number $101 \times 49$ is closest to: $50$; $500$; $5000$; $50\,000$? Then
compute it exactly with a quick method.
\guide
\step{1}{No brackets: multiply and divide first; add and subtract after}
\[ 20 - \underbrace{4\times 3}_{12} + \underbrace{8 \div 2}_{4} = 20 - 12 + 4 = 8 + 4 = 12. \]
(Subtraction and addition are done from left to right: $20 - 12 = 8$, then $8 + 4 = 12$.)
\step{2}{With brackets: brackets first}
\[ (20 - 4)\times(3+8) \div 2 = 16 \times 11 \div 2 = 176 \div 2 = 88. \]
\step{3}{Estimate}
$101 \approx 100$ and $49 \approx 50$, so $101 \times 49 \approx 100\times 50 = 5000$. Choose $5000$.
\step{4}{Exact value using the distributive property}
$101 \times 49 = (100 + 1)\times 49 = 4900 + 49 = 4949$, indeed very close to $5000$.
\ans{a) $12$ and $88$;\quad b) closest to $5000$; $101\times 49 = 4949$}
\end{exbox}

\begin{exbox}{Example 4}
a) Apples cost $\$1.80$ per kilogram. Grace buys $15$ kg and gets a $10\%$ discount on the whole bill. How much
does Grace pay?\\
b) A shirt costs $\$40$. After every week the shop reduces the \emph{current price} by $10\%$. After at least how
many weeks is the price of the shirt below $\$30$?
\guide
\step{1}{Total before the discount}
Total cost $=$ unit price $\times$ quantity $= 1.80 \times 15 = 27$ (dollars).
\step{2}{Subtract $10\%$}
$10\%$ of $27$ is $27 \div 10 = 2.70$ (dollars). Grace pays $27 - 2.70 = 24.30$ (dollars).
(Or: $90\%$ remains, i.e.\ $27 \times 0.9 = 24.3$.)
\step{3}{Successive discounts: multiply by $0.9$ every week}
Let $g_n$ be the price of the shirt (in dollars) after $n$ weeks. Then $g_0 = 40$ and $g_{n} = g_{n-1}\times 0.9$.
\begin{center}
\begin{tabular}{|c|c|c|c|c|}
\hline
\rowcolor{tmBlue!10} Time & Start & After 1 week & After 2 weeks & After 3 weeks\\ \hline
Price (dollars) & $40$ & $36$ & $32.40$ & $29.16$\\ \hline
Reduction (dollars) & -- & $4$ & $3.60$ & $3.24$\\ \hline
\end{tabular}
\end{center}
After $2$ weeks the price is $\$32.40 > \$30$; after $3$ weeks it is $\$29.16 < \$30$.
\begin{warnbox}
Do not think ``the price drops by $\$4$ every week'' (because $10\%$ of $40$ is $4$). Each week the reduction is
$10\%$ of \emph{the price at that time}, so the reductions get smaller and smaller: $4$, then $3.60$, then $3.24$, \ldots
\end{warnbox}
\ans{a) $\$24.30$;\quad b) $3$ weeks}
\end{exbox}

\begin{exbox}{Example 5}
a) Leo cycles $3$ laps around a stadium with a perimeter of $400$ m in $4$ minutes. Mia cycles at $15$ km/h. Who is faster?\\
b) A car travels at $72$ km/h. How many metres does it travel in $3$ seconds?
\guide
\step{1}{Leo's speed (m/s)}
Distance: $3 \times 400 = 1200$ (m). Time: $4$ minutes $= 4\times 60 = 240$ (seconds).
\[ v_{\text{Leo}} = 1200 \div 240 = 5 \text{ (m/s)}. \]
\step{2}{Convert Mia's speed to m/s}
\begin{recallbox}$1$ km/h $= \dfrac{1000 \text{ m}}{3600 \text{ s}} = \dfrac{5}{18}$ m/s.\end{recallbox}
$v_{\text{Mia}} = 15 \times \dfrac{5}{18} = \dfrac{75}{18} \approx 4.17$ (m/s). Since $5 > 4.17$, \textbf{Leo is faster}.
(Another way: $5$ m/s $= 5\times 3.6 = 18$ km/h $> 15$ km/h.)
\step{3}{The car}
$72$ km/h $= 72 \times \dfrac{5}{18} = 20$ m/s. In $3$ seconds it travels $20 \times 3 = 60$ (m).
\ans{a) Leo is faster ($5$ m/s against about $4.17$ m/s);\quad b) $60$ m}
\end{exbox}

\begin{exbox}{Example 6}
a) Which of the numbers $72;\ 84;\ 117;\ 135$ are divisible by $9$?\\
b) Find the digit $x$ such that the number $\nd{3x5}$ is divisible by $9$.\\
c) How many whole numbers are divisors of both $24$ and $36$?
\guide
\step{1}{Use the rule for $9$: compute the digit sum}
$72$: $7+2 = 9$ (divisible by $9$). \quad $84$: $8 + 4 = 12$ (not divisible by $9$).\\
$117$: $1+1+7 = 9$ (divisible). \quad $135$: $1+3+5 = 9$ (divisible).\\
So $72$, $117$, $135$ are divisible by $9$; $84$ is not ($84 = 9\times 9 + 3$).
\step{2}{Find the digit $x$}
The digit sum of $\nd{3x5}$ is $3 + x + 5 = 8 + x$. Since $x$ is a digit, $8 \le 8 + x \le 17$. In this range only $9$ is
divisible by $9$, so $8 + x = 9$, i.e.\ $x = 1$. The number is $315 = 9 \times 35$.
\step{3}{Common divisors of $24$ and $36$}
Divisors of $24$: $1, 2, 3, 4, 6, 8, 12, 24$. Divisors of $36$: $1, 2, 3, 4, 6, 9, 12, 18, 36$.\\
The numbers in both lists: $1, 2, 3, 4, 6, 12$ -- these are exactly the divisors of $\gcd(24,36) = 12$.
\ans{a) $72;\ 117;\ 135$;\quad b) $x = 1$;\quad c) $6$ numbers}
\end{exbox}

\begin{exbox}{Example 7}
In the calculation $\nd{AB} + \nd{AB} + \nd{AB} = \nd{1AB}$, each letter stands for a digit. Find $A$ and $B$.
\guide
\textbf{Method 1 -- Place value.} We have $\nd{1AB} = 100 + \nd{AB}$. So
\[ 3\times\nd{AB} = 100 + \nd{AB} \;\Rightarrow\; 2\times\nd{AB} = 100 \;\Rightarrow\; \nd{AB} = 50. \]
So $A = 5$, $B = 0$.

\textbf{Method 2 -- Column by column.} Ones column: $B + B + B = 3B$ ends in the digit $B$, so $3B - B = 2B$ ends in $0$,
hence $B = 0$ or $B = 5$.
\begin{itemize}[leftmargin=1.4em,itemsep=1pt,topsep=1pt]
  \item If $B = 5$: $3\times 5 = 15$, write $5$ carry $1$. Tens column: $3A + 1 = \nd{1A} = 10 + A$, so $2A = 9$ -- no digit
  works. Rejected.
  \item If $B = 0$: no carry. Tens column: $3A = 10 + A$, so $2A = 10$, $A = 5$.
\end{itemize}
\verify{$50 + 50 + 50 = 150 = \nd{1AB}$ with $A = 5$, $B = 0$. Correct.}
\ans{$A = 5$, $B = 0$}
\end{exbox}

\begin{exbox}{Example 8}
a) Noah's four test scores are $6$; $9$; $10$ and $x$. The mean of the four scores is $8$. Find $x$.\\
b) The mean of five consecutive whole numbers is $30$. Find the smallest of them.
\guide
\step{1}{Form an equation from the definition of the mean}
\begin{recallbox}Mean $=$ sum $\div$ how many numbers, so sum $=$ mean $\times$ how many numbers.\end{recallbox}
\[ \frac{6 + 9 + 10 + x}{4} = 8 \;\Rightarrow\; 25 + x = 32 \;\Rightarrow\; x = 32 - 25 = 7. \]
\step{2}{Using deviations (faster)}
Compared with $8$: the $6$ is $2$ short, the $9$ is $1$ over, the $10$ is $2$ over. Sum of deviations: $-2 + 1 + 2 = +1$.
For all deviations to add up to $0$, the fourth score must be $1$ short: $x = 8 - 1 = 7$.
\step{3}{Five consecutive numbers}
Let the smallest number be $n$; the five numbers are $n, n+1, n+2, n+3, n+4$, with sum $5n + 10$. So
$5n + 10 = 5\times 30 = 150$, hence $5n = 140$, $n = 28$. (Note: the middle number $n + 2 = 30$ is exactly the mean.)
The five numbers are $28, 29, 30, 31, 32$; the smallest is $28$.
\verify{$28 + 29 + 30 + 31 + 32 = 150 = 5\times 30$. Correct.}
\ans{a) $x = 7$;\quad b) $28$}
\end{exbox}

% =====================================================================
\clearpage
\dividerbar{Core practice}

% ---------------------------------------------------------------- Problem 1
\needspace{10\baselineskip}
\begin{probbox}{1}{\relax}
The smallest number in the set $\{3.2;\ 2.3;\ 3;\ 2.23;\ 3.22\}$ is
\choices{$3.2$}{$2.3$}{$3$}{$2.23$}{$3.22$}
\end{probbox}
\probsol{
\textbf{Understanding the problem.} There are $5$ decimals. We need the \emph{smallest} one.
\begin{recallbox}Comparing decimals: compare the whole parts first; if they are equal, compare the tenths, then the
hundredths, \ldots\ You may add zeros at the end of the decimal part without changing the value.\end{recallbox}
\step{1}{Compare the whole parts}
Whole parts: $3.2 \to 3$; \ $2.3 \to 2$; \ $3 \to 3$; \ $2.23 \to 2$; \ $3.22 \to 3$.\\
The smallest number must have the smallest whole part, namely $2$. Only two candidates remain: $2.3$ and $2.23$.
\step{2}{Use the same number of decimal places and compare}
$2.3 = 2.30$ \ and \ $2.23$. Tenths: $3$ against $2$; since $2 < 3$, we get $2.23 < 2.30$.
\step{3}{Conclusion}
Ordering everything to be sure: $2.23 < 2.3 < 3 < 3.2 < 3.22$. The smallest number is $2.23$.
\begin{warnbox}Do not choose $2.3$ just because ``$3$ is less than $23$''. The digit $3$ in $2.3$ means $3$ tenths, while
$2.23$ has only $2$ tenths.\end{warnbox}
\verify{$2.23 = \dfrac{223}{100}$ and $2.3 = \dfrac{230}{100}$; $223 < 230$.}
The number $2.23$ is option \textbf{(D)}.
\ans{(D) $2.23$}
}

% ---------------------------------------------------------------- Problem 2
\needspace{10\baselineskip}
\begin{probbox}{2}{\relax}
Keegan goes on a $36$ km kayak trip. Before lunch he paddles the first $12$ km. After lunch, what fraction of the whole
trip does Keegan still have to paddle?
\choices{$\dfrac{1}{2}$}{$\dfrac{5}{6}$}{$\dfrac{3}{4}$}{$\dfrac{2}{3}$}{$\dfrac{3}{5}$}
\end{probbox}
\probsol{
\textbf{Understanding the problem.} The whole trip is $36$ km and $12$ km are done. We are asked for the part \emph{left}
(not yet paddled), as a fraction.
\step{1}{Distance left}
Let $x$ (km) be the distance still to paddle after lunch. Then $12 + x = 36$, so $x = 36 - 12 = 24$ (km).
\step{2}{Form the fraction ``part left over the whole''}
\begin{recallbox}The fraction ``what part of the whole'' $=\dfrac{\text{the part}}{\text{the whole}}$. Simplify by dividing
numerator and denominator by their greatest common divisor.\end{recallbox}
\[ \frac{24}{36} = \frac{24 \div 12}{36 \div 12} = \frac{2}{3}. \]
\step{3}{Another way: find the part done and subtract it from $1$}
Part done: $\dfrac{12}{36} = \dfrac{1}{3}$. Part left: $1 - \dfrac13 = \dfrac33 - \dfrac13 = \dfrac23$. Both ways give the same result.
\begin{warnbox}The fraction $\dfrac{12}{36} = \dfrac{1}{3}$ is the part \emph{done}. The question asks for the part
\emph{left}, so the answer is $1 - \dfrac13 = \dfrac23$.\end{warnbox}
\verify{$\dfrac{2}{3}\times 36 = 24$ km; $24 + 12 = 36$ km. Correct.}
\begin{parentbox}Ask your child two questions: ``What part is done?'' and ``What part is left?''. Getting into the habit of
answering both helps avoid the trap of choosing $\dfrac13$.\end{parentbox}
The fraction $\dfrac{2}{3}$ is option \textbf{(D)}.
\ans{(D) $\dfrac{2}{3}$}
}

% ---------------------------------------------------------------- Problem 3
\needspace{10\baselineskip}
\begin{probbox}{3}{\relax}
Which of the following statements is \textbf{true}?
\choicesII{$0$ is less than $-5$}{$7$ is less than $-1$}{$10$ is less than $\dfrac{1}{4}$}{$-1$ is less than $-3$}{$-8$ is less than $-2$}
\end{probbox}
\probsol{
\textbf{Understanding the problem.} Each option is a comparison. We check them one by one to find the true one.
\begin{recallbox}On a horizontal number line, a number to the left is smaller than a number to the right. Every negative
number is less than $0$ and less than every positive number. Of two negative numbers, the one farther from $0$ is smaller.\end{recallbox}
\step{1}{Place the numbers on a number line}
\begin{figC}
  \draw[->,thick,Navy] (-9.4*0.6,0)--(10.6*0.6,0);
  \foreach \x in {-9,...,10}{\draw[Navy] (\x*0.6,0.08)--(\x*0.6,-0.08);}
  \foreach \x in {-8,-5,-3,-2,-1,0,7,10}{\node[below,font=\scriptsize] at (\x*0.6,-0.1) {$\x$};}
  \foreach \x in {-8,-2}{\fill[tmOrange] (\x*0.6,0) circle (3pt);}
  \draw[->,tmGreen!70!black,thick,shorten <=3.5pt,shorten >=3.5pt] (-8*0.6,0) to[bend left=50] node[above,font=\scriptsize]{$6$ units to the right} (-2*0.6,0);
\end{figC}
\step{2}{Check each option}
(A) $0$ is less than $-5$? False: $-5$ is to the left of $0$, so $-5 < 0$.\\
(B) $7$ is less than $-1$? False: a positive number is always greater than a negative number.\\
(C) $10$ is less than $\dfrac14$? False: $\dfrac14 = 0.25 < 10$.\\
(D) $-1$ is less than $-3$? False: $-3$ is farther from $0$, so $-3 < -1$.\\
(E) $-8$ is less than $-2$? \textbf{True}: $-8$ is to the left of $-2$ (farther from $0$).
\begin{warnbox}Many students choose (D), thinking ``$1 < 3$ so $-1 < -3$''. For negative numbers it is the other way round:
$-1 > -3$ (a temperature of $-1^\circ$C is warmer than $-3^\circ$C).\end{warnbox}
\verify{$-2 - (-8) = 6 > 0$, which means $-2$ is greater than $-8$.}
The only true statement is \textbf{(E)}.
\ans{(E) $-8$ is less than $-2$}
}

% ---------------------------------------------------------------- Problem 4
\needspace{10\baselineskip}
\begin{probbox}{4}{\relax}
The value of $6 \times 2 - 3$ is
\choices{$9$}{$-6$}{$12$}{$15$}{$10$}
\end{probbox}
\probsol{
\textbf{Understanding the problem.} The expression has two operations, multiplication and subtraction, and no brackets.
\begin{recallbox}In an expression without brackets: multiply and divide first; add and subtract after.\end{recallbox}
\step{1}{Do the multiplication first}
$6\times 2 = 12$.
\step{2}{Then the subtraction}
$12 - 3 = 9$.
\begin{warnbox}Subtracting first would give $6\times(2-3) = 6\times(-1) = -6$ -- option (B), a ``trap'' for anyone who
forgets the order of operations. The expression has \emph{no brackets}, so we must multiply first.\end{warnbox}
\verify{Read the expression as a story: ``$6$ boxes with $2$ cakes each, $3$ cakes eaten'' leaves $12 - 3 = 9$ cakes.}
The value $9$ is option \textbf{(A)}.
\ans{(A) $9$}
}

% ---------------------------------------------------------------- Problem 5
\needspace{10\baselineskip}
\begin{probbox}{5}{\relax}
A test is marked out of $50$. Mukesh scores $80\%$. How many marks does Mukesh get?
\choices{$40$}{$62.5$}{$10$}{$45$}{$35$}
\end{probbox}
\probsol{
\textbf{Understanding the problem.} The maximum is $50$ marks; Mukesh gets $80\%$ of $50$. We need Mukesh's marks.
\begin{recallbox}To find $p\%$ of a number $a$, compute $a \times \dfrac{p}{100}$.\end{recallbox}
\step{1}{Write the calculation}
$80\%$ of $50$ is $50 \times \dfrac{80}{100}$.
\step{2}{Calculate}
$50 \times \dfrac{80}{100} = \dfrac{50\times 80}{100} = \dfrac{4000}{100} = 40$.
\step{3}{A quick mental method}
$10\%$ of $50$ is $5$; so $80\% = 8\times 10\%$ is $8\times 5 = 40$.
\begin{warnbox}Option (B) $62.5$ comes from dividing $50 \div 0.8$ -- wrong, since $62.5$ is even more than the maximum $50$!
Option (C) $10$ is the number of marks Mukesh \emph{lost} ($20\%$ of $50$).\end{warnbox}
\verify{$\dfrac{40}{50} = \dfrac{80}{100} = 80\%$. Correct.}
The value $40$ is option \textbf{(A)}.
\ans{(A) $40$}
}

% ---------------------------------------------------------------- Problem 6
\needspace{10\baselineskip}
\begin{probbox}{6}{\relax}
Ground beef sells for $\$5.00$ per kilogram. How much do $12$ kg of ground beef cost?
\choices{$\$5.00$}{$\$12.00$}{$\$60.00$}{$\$17.00$}{$\$2.40$}
\end{probbox}
\probsol{
\textbf{Understanding the problem.} The unit price (price of $1$ kg) is $\$5.00$; the quantity is $12$ kg. We need the total cost.
\begin{recallbox}Total cost $=$ unit price $\times$ quantity.\end{recallbox}
\step{1}{Write the calculation}
Let $T$ (dollars) be the amount to pay. Then $T = 5 \times 12$.
\step{2}{Calculate}
$T = 5\times 12 = 60$ (dollars), i.e.\ $\$60.00$.
\step{3}{Estimate to check}
Each kilogram costs $\$5$ and we buy more than $10$ kg, so the cost is more than $5\times 10 = 50$ dollars. Only option (C) fits.
\begin{warnbox}The wrong options come from using the wrong operation: (A) is the price of $1$ kg; (B) is just the number of
kilograms written as money; (D) is $5 + 12$ (adding instead of multiplying); (E) is $12 \div 5$ (dividing instead of multiplying).\end{warnbox}
\verify{$60 \div 12 = 5$ dollars per kilogram. Correct.}
The amount $\$60.00$ is option \textbf{(C)}.
\ans{(C) $\$60.00$}
}

% ---------------------------------------------------------------- Problem 7
\needspace{10\baselineskip}
\begin{probbox}{7}{\relax}
The value of $(8 \times 4) + 3$ is
\choices{$96$}{$15$}{$56$}{$35$}{$28$}
\end{probbox}
\probsol{
\textbf{Understanding the problem.} The brackets contain a multiplication; then $3$ is added.
\begin{recallbox}Expressions with brackets: work out the brackets first.\end{recallbox}
\step{1}{Work out the brackets}
$8\times 4 = 32$.
\step{2}{Add $3$}
$32 + 3 = 35$.
\begin{warnbox}Misreading it as $8 \times (4 + 3) = 8\times 7 = 56$ leads to option (C). The brackets surround $8\times 4$,
not $4 + 3$. Option (A) $96 = 8\times 4\times 3$ comes from mistaking the plus sign for a times sign.\end{warnbox}
\verify{Without brackets the result is still $8\times 4 + 3 = 35$, because multiplication always comes before addition; here the
brackets only add emphasis.}
The value $35$ is option \textbf{(D)}.
\ans{(D) $35$}
}

% ---------------------------------------------------------------- Problem 8
\needspace{10\baselineskip}
\begin{probbox}{8}{\relax}
Which of the following integers is closest to $99 \times 9$?
\choices{$10\,000$}{$100$}{$100\,000$}{$1000$}{$10$}
\end{probbox}
\probsol{
\textbf{Understanding the problem.} The options are $10$; $100$; $1000$; $10\,000$; $100\,000$ -- each $10$ times the
previous. An estimate of $99\times 9$ is enough, but we also compute it exactly to be sure.
\step{1}{Estimate}
$99 \approx 100$, so $99 \times 9 \approx 100\times 9 = 900$.
\step{2}{Exact value using the distributive property}
\begin{recallbox}Distributive property: $(a - b)\times c = a\times c - b\times c$.\end{recallbox}
$99\times 9 = (100 - 1)\times 9 = 900 - 9 = 891$.
\step{3}{Compare the distances to the nearest options}
The distance from $891$ to $1000$ is $1000 - 891 = 109$.\\
The distance from $891$ to $100$ is $891 - 100 = 791$.\\
The other options ($10$; $10\,000$; $100\,000$) are much farther away. So $1000$ is the closest.
\begin{figC}
  \draw[->,thick,Navy] (-0.3,0)--(10.6,0);
  \foreach \x/\l in {0/0,1/100,10/1000}{\draw[Navy] (\x,0.1)--(\x,-0.1); \node[below,font=\scriptsize] at (\x,-0.1) {$\l$};}
  \fill[tmOrange] (8.91,0) circle (3pt); \node[above,font=\scriptsize] at (8.91,0.1) {$891$};
  \draw[Slate!60,densely dotted] (1,-0.5)--(1,-0.95) (8.91,-0.1)--(8.91,-0.95) (10,-0.5)--(10,-0.95);
  \draw[<->,tmGreen!70!black] (8.91,-0.85)--(10,-0.85) node[midway,below,font=\scriptsize]{$109$};
  \draw[<->,AnsRed!70] (1,-0.85)--(8.91,-0.85) node[midway,below,font=\scriptsize]{$791$};
\end{figC}
\verify{$891 \div 9 = 99$. Correct.}
The number $1000$ is option \textbf{(D)}.
\ans{(D) $1000$}
}

% ---------------------------------------------------------------- Problem 9
\needspace{10\baselineskip}
\begin{probbox}{9}{\relax}
Joshua is reading a $396$-page book. He has read only the first $\dfrac{1}{3}$ of the book. How many pages does he still
have to read to finish the book?
\choices{$264$}{$124$}{$250$}{$199$}{$244$}
\end{probbox}
\probsol{
\textbf{Understanding the problem.} The whole book: $396$ pages. Read so far: $\dfrac13$ of the book. Asked: the number of
pages \emph{left}.
\begin{recallbox}To find $\dfrac{m}{n}$ of a number $a$, compute $a\times\dfrac{m}{n}$ (divide $a$ by $n$ and multiply by $m$).\end{recallbox}
\step{1}{Pages read}
$\dfrac13$ of $396$ is $396\times\dfrac13 = 396 \div 3 = 132$ (pages).
\step{2}{Introduce an unknown and form an equation}
Let $x$ be the number of pages Joshua still has to read. Pages read plus pages left make the whole book:
\[ 132 + x = 396 \;\Rightarrow\; x = 396 - 132 = 264. \]
\step{3}{Another way: compute the part left directly}
The part left is $1 - \dfrac13 = \dfrac23$ of the book, i.e.\ $396 \times \dfrac23 = \dfrac{396\times 2}{3} = \dfrac{792}{3} = 264$ (pages).
\begin{warnbox}Stopping at $132$ answers the wrong question (that is the number of pages \emph{read}; and $132$ is not even
among the options). Always re-read the question before choosing an answer.\end{warnbox}
\verify{$264 + 132 = 396$, and $132\times 3 = 396$. Correct.}
The number $264$ is option \textbf{(A)}.
\ans{(A) $264$}
}

% ---------------------------------------------------------------- Problem 10
\needspace{10\baselineskip}
\begin{probbox}{10}{\relax}
Which of the following numbers is a multiple of $9$?
\choices{$50$}{$40$}{$35$}{$45$}{$55$}
\end{probbox}
\probsol{
\textbf{Understanding the problem.} We need the number divisible by $9$ (a multiple of $9$) among the five given numbers.
\begin{recallbox}A number is divisible by $9$ exactly when its digit sum is divisible by $9$.\end{recallbox}
\step{1}{Digit sum of each number}
\begin{center}
\begin{tabular}{|c|c|c|c|c|c|}
\hline
\rowcolor{tmBlue!10} Number & $50$ & $40$ & $35$ & $45$ & $55$\\ \hline
Digit sum & $5$ & $4$ & $8$ & $9$ & $10$\\ \hline
Divisible by $9$? & no & no & no & \textbf{yes} & no\\ \hline
\end{tabular}
\end{center}
\step{2}{Check with the $9$ times table}
Multiples of $9$ nearby: $36,\ 45,\ 54$. Only $45 = 9\times 5$ appears among the options.
\begin{proofbox}Why does the rule work for a two-digit number $\nd{ab}$? We have $\nd{ab} = 10a + b = 9a + (a + b)$. The part $9a$
is divisible by $9$, so $\nd{ab}$ is divisible by $9$ exactly when $a + b$ is. For $45$: $45 = 9\times 4 + 9$.\end{proofbox}
\verify{$45 \div 9 = 5$, remainder $0$.}
The number $45$ is option \textbf{(D)}.
\ans{(D) $45$}
}

% ---------------------------------------------------------------- Problem 11
\needspace{10\baselineskip}
\begin{probbox}{11}{\relax}
Consider the set $\{0.34;\ 0.304;\ 0.034;\ 0.43\}$. The sum of the smallest and the largest numbers in the set is
\choices{$0.77$}{$0.734$}{$0.077$}{$0.464$}{$0.338$}
\end{probbox}
\probsol{
\textbf{Understanding the problem.} Part 1: find the smallest and the largest number. Part 2: add them.
\begin{recallbox}To compare decimals, add zeros so the decimal parts have equal length, then compare place by place from
left to right.\end{recallbox}
\step{1}{Write the numbers with three decimal places}
\begin{center}
\begin{tabular}{|c|c|c|c|}
\hline
\rowcolor{tmBlue!10} Original & Rewritten & Tenths & Hundredths \\ \hline
$0.34$ & $0.340$ & $3$ & $4$ \\ \hline
$0.304$ & $0.304$ & $3$ & $0$ \\ \hline
$0.034$ & $0.034$ & $0$ & $3$ \\ \hline
$0.43$ & $0.430$ & $4$ & $3$ \\ \hline
\end{tabular}
\end{center}
\step{2}{Find the smallest and the largest}
Comparing like whole numbers: $34 < 304 < 340 < 430$, so $0.034 < 0.304 < 0.34 < 0.43$.\\
Smallest: $0.034$ (its tenths digit is $0$). Largest: $0.43$ (its tenths digit is $4$).
\step{3}{Add, with the decimal points lined up}
\[ \begin{array}{r} 0.430\\ +\ 0.034\\ \hline 0.464 \end{array} \]
\begin{warnbox}Misaligning the columns (adding $43$ and $34$ ``by the last digits'') gives $0.77$ or $0.077$ -- options (A)
and (C). Keep the decimal points in line: $4$ tenths plus $0$ tenths, $3$ hundredths plus $3$ hundredths, \ldots\end{warnbox}
\verify{$0.464 - 0.43 = 0.034$. Correct.}
The sum $0.464$ is option \textbf{(D)}.
\ans{(D) $0.464$}
}

% ---------------------------------------------------------------- Problem 12
\needspace{11\baselineskip}
\begin{probbox}{12}{\relax}
Keiko and Leah run on a circular track $150$ m long. Keiko runs $3$ laps in $120$ seconds, and Leah runs $5$ laps in $160$
seconds. Who is the faster runner, and approximately what is her speed?
\choicesII{Keiko, $3.75$ m/s}{Keiko, $2.4$ m/s}{Leah, $3.3$ m/s}{Leah, $4.69$ m/s}{Leah, $3.75$ m/s}
\end{probbox}
\probsol{
\textbf{Understanding the problem.} One lap is $150$ m. Keiko: $3$ laps in $120$ s. Leah: $5$ laps in $160$ s. We compute
each girl's speed (m/s), compare, and choose the option.
\begin{recallbox}Speed $=$ distance $\div$ time. The distance of $n$ laps $= n\times$ the length of one lap.\end{recallbox}
\step{1}{Distance run by each girl}
Keiko: $3\times 150 = 450$ (m). \quad Leah: $5\times 150 = 750$ (m).
\step{2}{Keiko's speed}
$v_{\text{Keiko}} = 450 \div 120 = 3.75$ (m/s).
\step{3}{Leah's speed}
$v_{\text{Leah}} = 750 \div 160 = 4.6875 \approx 4.69$ (m/s).
\step{4}{Compare}
$4.69 > 3.75$, so \textbf{Leah is faster}, running at about $4.69$ m/s.
\begin{center}
\begin{tabular}{|c|c|c|c|}
\hline
\rowcolor{tmBlue!10} & Distance & Time & Speed\\ \hline
Keiko & $450$ m & $120$ s & $3.75$ m/s\\ \hline
Leah & $750$ m & $160$ s & $\approx 4.69$ m/s\\ \hline
\end{tabular}
\end{center}
\begin{warnbox}Option (E) ``Leah, $3.75$ m/s'' pairs the right name with the slower girl's speed. Option (B) $2.4$ m/s comes from
the reversed calculation $120 \div 50$ (time divided by distance). Always write the units: m $\div$ s $=$ m/s.\end{warnbox}
\verify{$4.6875 \times 160 = 750$ m $= 5$ laps; $3.75\times 120 = 450$ m $= 3$ laps. Correct.}
The conclusion ``Leah, $4.69$ m/s'' is option \textbf{(D)}.
\ans{(D) Leah, $4.69$ m/s}
}

% ---------------------------------------------------------------- Problem 13
\needspace{16\baselineskip}
\begin{probbox}{13}{\relax}
In the addition shown, $P$ and $Q$ each represent a digit. The value of $P + Q$ is
\[ \begin{array}{r@{\;}c@{\;}c@{\;}c} & P & Q & Q\\ & P & P & Q\\ + & Q & Q & Q\\ \hline & 8 & 7 & 6 \end{array} \]
\choices{$3$}{$5$}{$7$}{$6$}{$4$}
\end{probbox}
\probsol{
\textbf{Understanding the problem.} The three numbers $\nd{PQQ}$, $\nd{PPQ}$, $\nd{QQQ}$ add up to $876$. We need the digits $P$, $Q$,
then $P + Q$. We go column by column, starting with the ones.
\step{1}{Ones column: $Q + Q + Q = 3Q$ ends in $6$}
Since $Q$ is a digit, $3Q$ is one of $0, 3, 6, \ldots, 27$. Of these, only $6$ ends in $6$ ($16$ and $26$ are not multiples
of $3$). So $3Q = 6$, $Q = 2$, and there is \emph{no carry} into the tens.
\step{2}{Tens column: $Q + P + Q = P + 4$ ends in $7$}
$P + 4$ lies between $4$ and $13$; ending in $7$ means $P + 4 = 7$, so $P = 3$. No carry into the hundreds.
\step{3}{Hundreds column: check $P + P + Q = 8$}
$3 + 3 + 2 = 8$. This matches the hundreds digit of $876$.
\begin{recallbox}Place value: $\nd{PQQ} = 100P + 11Q$, \ $\nd{PPQ} = 110P + Q$, \ $\nd{QQQ} = 111Q$.\end{recallbox}
\step{4}{Another way using place value}
Adding: $(100P + 11Q) + (110P + Q) + 111Q = 210P + 123Q = 876$. Try $Q = 2$: $210P = 876 - 246 = 630$, so $P = 3$.
(Other values of $Q$ do not give a digit for $P$.)
\verify{$322 + 332 + 222 = 876$. Correct.}
\begin{parentbox}A digit puzzle is like a detective game: each column is a ``clue''. Let your child write each column as a
sentence (``three times $Q$ ends in $6$'') and try the digits $0$ to $9$. Do not forget to check the carries.\end{parentbox}
So $P + Q = 3 + 2 = 5$, which is option \textbf{(B)}.
\ans{(B) $5$}
}

% ---------------------------------------------------------------- Problem 14
\needspace{12\baselineskip}
\begin{probbox}{14}{\relax}
The mean (average) of the four integers $78$; $83$; $82$ and $x$ is $80$. Which of the following statements is true?
\choicesL{$x$ is $2$ greater than the mean}{$x$ is $1$ less than the mean}{$x$ is $2$ less than the mean}{$x$ is $3$ less than the mean}{$x$ is equal to the mean}
\end{probbox}
\probsol{
\textbf{Understanding the problem.} We know the mean of four numbers is $80$ and we know three of them. Find $x$ and compare it with $80$.
\begin{recallbox}Sum of the numbers $=$ mean $\times$ how many numbers.\end{recallbox}
\step{1}{Form an equation}
By the definition of the mean: $\dfrac{78 + 83 + 82 + x}{4} = 80$, i.e.\ $78 + 83 + 82 + x = 80\times 4 = 320$.
\step{2}{Solve the equation for $x$}
$78 + 83 + 82 = 243$, so $243 + x = 320$, hence $x = 320 - 243 = 77$.
\step{3}{Compare with the mean}
$80 - 77 = 3$: $x$ is $3$ less than the mean.
\step{4}{A faster way -- the deviation method}
Compared with $80$: $78$ is $2$ short; $83$ is $3$ over; $82$ is $2$ over. Sum of these three deviations: $-2 + 3 + 2 = +3$.
Since all the deviations must add up to $0$, $x$ must be $3$ ``short'': $x = 80 - 3 = 77$.
\begin{proofbox}Each number can be written as $80 + $(its deviation). The sum of the four numbers is $4\times 80 + $(sum of deviations).
But the sum of the four numbers is exactly $4\times 80$, so the sum of the deviations is $0$.\end{proofbox}
\verify{$(78 + 83 + 82 + 77) \div 4 = 320 \div 4 = 80$. Correct.}
\begin{parentbox}The ``deviation'' method is great for working out average marks mentally in everyday life. You can ask your
child: ``Three tests scored $7$, $9$, $8$; to average $8$, what does the fourth test need?'' (Answer: $8$, because
$-1 + 1 + 0 = 0$.)\end{parentbox}
The statement ``$x$ is $3$ less than the mean'' is option \textbf{(D)}.
\ans{(D)}
}

% ---------------------------------------------------------------- Problem 15
\needspace{10\baselineskip}
\begin{probbox}{15}{\relax}
In the Fibonacci sequence $1, 1, 2, 3, 5, \ldots$, each number from the $2$ onwards is the sum of the two numbers just
before it. For example, the next number in the sequence is $3 + 5 = 8$. Which of the following numbers is in the sequence?
\choices{$20$}{$21$}{$22$}{$23$}{$24$}
\end{probbox}
\probsol{
\textbf{Understanding the problem.} Rule: next number $=$ sum of the two before it. We continue the sequence until it
passes $24$ (the largest option), then see which option appears.
\step{1}{Continue the sequence}
\begin{center}
\begin{tabular}{|c|c|c|c|c|c|c|c|c|c|}
\hline
\rowcolor{tmBlue!10} Position & 1 & 2 & 3 & 4 & 5 & 6 & 7 & 8 & 9\\ \hline
Number & $1$ & $1$ & $2$ & $3$ & $5$ & $8$ & $13$ & $\mathbf{21}$ & $34$\\ \hline
\end{tabular}
\end{center}
In detail: $5 + 8 = 13$; \ $8 + 13 = 21$; \ $13 + 21 = 34$.
\step{2}{Compare with the options}
$34$ is already greater than $24$, so all later numbers are greater than $24$. Between $13$ and $34$ there is exactly one
number of the sequence, $21$. So $20$, $22$, $23$, $24$ are not in the sequence; only $21$ is.
\begin{tipbox}The Fibonacci sequence keeps increasing (from the second number on) because we keep adding positive numbers.
So once the sequence has passed a number, that number can never appear later.\end{tipbox}
\verify{$8 + 13 = 21$. Correct.}
The number $21$ is option \textbf{(B)}.
\ans{(B) $21$}
}
